\documentclass[10pt,a4paper]{article}


%
%

%
%
%

\usepackage[utf8]{inputenc}
%
\usepackage{xspace}
\usepackage{amsmath}
\usepackage{amsthm}
\usepackage{amsfonts}
\usepackage{amssymb}
\usepackage{mathtools}
\usepackage{textcomp}
\usepackage{varwidth}
\usepackage{graphicx}
\usepackage{mathrsfs} 

\usepackage{hyperref}
\usepackage{caption}
\usepackage{subcaption}
\usepackage{enumerate}
\usepackage{xcolor}
\usepackage[noadjust]{cite}

\usepackage{tikz}
\usetikzlibrary{automata, graphs,positioning,chains,arrows,calc, backgrounds}

%

%
%
%



\newcommand{\defining}[1]{\emph{#1}}

\DeclarePairedDelimiter\set{\lbrace}{\rbrace}

\DeclarePairedDelimiterX\setcond[2]{\{}{\}}{\mathchoice{\,}{}{}{}#1 \;\delimsize\vert\; #2\mathchoice{\,}{}{}{}}


\newcommand{\regular}{\mathsf{regular}}
\newcommand{\blocking}{\mathsf{blocking}}

%
\renewcommand{\phi}{\varphi}


\newcommand{\inv}[1]{#1^{-1}}

\newcommand{\Vars}[1]{\mathsf{var}(#1)}

\newcommand{\comp}{\equiv}
\newcommand{\CFIsym}{\mathsf{CFI}}
\newcommand{\CFI}[1]{\CFIsym(#1)}
\newcommand{\compCFIsym}{\mathsf{cCFI}}
\newcommand{\compCFI}[2]{\CFI{#1}/_{#2}}
\newcommand{\precompCFI}[2]{{(\CFI{#1},#2)}}

\newcommand{\nat}{\mathbb{N}}
\newcommand{\FF}{\mathbb{F}}
\newcommand{\tup}[1]{\bar{#1}}
\newcommand{\iso}{\cong}

%
%
%
\newcommand{\klogic}[1]{\mathscr{L}_{#1}}
\newcommand{\kmlogic}[2]{\mathscr{L}_{#1, #2}}

\newcommand{\graphG}{\mathcal{G}}
\newcommand{\graphH}{\mathcal{H}}
\newcommand{\grid}{\mathcal{C}}
\newcommand{\vertexG}{V_{\graphG}}
\newcommand{\vertexH}{V_{\graphH}}
\newcommand{\vertexC}{V_{\mathcal{C}}}
\newcommand{\edgeG}{E_{\graphG}}
\newcommand{\edgeH}{E_{\graphH}}
\newcommand{\edgeC}{E_{\mathcal{C}}}
\newcommand{\coloringG}{\chi_{\graphG}}
\newcommand{\coloringH}{\chi_{\graphH}}
\newcommand{\labelG}{\gamma_{\graphG}}
\newcommand{\labelH}{\gamma_{\graphH}}

\newcommand{\var}{\mathsf{var}}
\newcommand{\dom}{\mathsf{dom}}

\newcommand{\graphXG}{\mathcal{X}(\mathcal{G})}
\newcommand{\graphXH}{\mathcal{X}(\mathcal{H})}
\newcommand{\vertexXG}{V_{\graphXG}}
\newcommand{\vertexXH}{V_{\graphXH}}
\newcommand{\edgeXG}{E_{\graphXG}}
\newcommand{\edgeXH}{E_{\graphXH}}
\newcommand{\labelXG}{\gamma_{\graphXG}}
\newcommand{\labelXH}{\gamma_{\graphXH}}

\newcommand{\ISO}{\mathsf{ISO}(\graphG, \graphH)}
\newcommand{\XISO}{\mathsf{ISO}(\graphXG, \graphXH)}


\newcommand{\kbequivr}[2]{\simeq_{\mathcal{B}}^{#1,#2}}
\newcommand{\notkbequivr}[2]{\not\simeq_{\mathcal{B}}^{#1,#2}}
\newcommand{\kbequiv}[1]{\simeq_{\mathcal{B}}^{#1}}
\newcommand{\kequiv}[1]{\simeq^{#1}}
\newcommand{\notkequiv}[1]{\not\simeq^{#1}}
\newcommand{\kequivr}[2]{\simeq^{#1,#2}}
\newcommand{\notkequivr}[2]{\not\simeq^{#1,#2}}


\newcommand{\X}{\mathcal{X}}
\newcommand{\XG}{\X_{\graphG}}
\newcommand{\XH}{\X_{\graphH}}
\newcommand{\XGI}{\X^{-1}_{\graphG}}
\newcommand{\XHI}{\X^{-1}_{\graphH}}
\newcommand{\XI}{\X^{-1}}


%
%


\definecolor{rwth}   {RGB}{  0  84 159}
\definecolor{rwth-75}{RGB}{ 64 127 183}
\definecolor{rwth-50}{RGB}{142 186 229}
\definecolor{rwth-25}{RGB}{199 221 242}
\definecolor{rwth-10}{RGB}{232 241 250}

\definecolor{darkblue}{rgb}{0,0,0.8}
\definecolor{darkgreen}{rgb}{0,0.7,0}
\definecolor{darkred}{rgb}{0.5,0,0}
\definecolor{darkyellow}{rgb}{0.6,0.6,0}

\definecolor{mediumblue}{rgb}{0.2,0.2,1}
\definecolor{lightblue}{rgb}{0.5,0.5,1}
\definecolor{lightgreen}{rgb}{0.5,1,0.5}
\definecolor{lightred}{rgb}{1,0.5,0.5}
\definecolor{lightyellow}{rgb}{1,1,0.5}


\definecolor{colA}{RGB}{0,80,250}
\definecolor{colB}{RGB}{60,200,230}
\definecolor{colC}{RGB}{60,180,75}
\definecolor{colD}{RGB}{230,25,75}
\definecolor{colE}{RGB}{245,130,48}
\definecolor{colF}{RGB}{145,30,180}
\definecolor{colG}{RGB}{240,50,230}
\definecolor{colH}{RGB}{170,110,40}



\tikzstyle{vertex} = [circle, fill=black, inner sep=0mm, minimum size = 2mm]

\tikzstyle{basevertex} = [circle, fill=black, inner sep=0mm, minimum size = 6mm]

\newcommand{\outerVertexDist}{0.3}
\newcommand{\gadgetRadius}{0.5}
\newcommand{\gadgetVertexDist}{0.08}
\newcommand{\gadgetScale}{0.8}


\newcommand{\drawEdgeVertexPair}[2]{
        \begin{scope}[#2] 
                \begin{scope}[on background layer]
                        \draw[draw, very thin, gray] (\gadgetRadius,0) ellipse (0.15cm and 0.32cm);
                \end{scope}
                \node[edgeVertex] (#1-0) at (\gadgetRadius,-\outerVertexDist/2) {};
                \node[edgeVertex] (#1-1) at (\gadgetRadius, \outerVertexDist/2) {};
        \end{scope}
}

\newcommand{\drawGadgetVertices}{
        \node[edgeVertex] (gul) at (-\gadgetVertexDist,\gadgetVertexDist) {};
        \node[edgeVertex] (gur) at (\gadgetVertexDist,\gadgetVertexDist) {};
        \node[edgeVertex] (gdl) at (-\gadgetVertexDist,-\gadgetVertexDist) {};
        \node[edgeVertex] (gdr) at (\gadgetVertexDist,-\gadgetVertexDist) {};
}

\newcommand{\drawGadgetCenter}[2]{
        \begin{scope}[name prefix=#1, shift={#2}, scale=\gadgetScale]
                \drawEdgeVertexPair{vu}{rotate=90}
                \drawEdgeVertexPair{vr}{rotate=0}
                \drawEdgeVertexPair{vd}{rotate=-90}
                \drawEdgeVertexPair{vl}{rotate=-180}
                \drawGadgetVertices{}
                
                \draw[black, draw]
                (vu-0) -- (gul)
                (vu-0) -- (gdr)
                (vu-1) -- (gur)
                (vu-1) -- (gdl);
                
                \draw[black, draw]
                (vr-0) -- (gur)
                (vr-0) -- (gdl)
                (vr-1) -- (gul)
                (vr-1) -- (gdr);
                
                \draw[black, draw]
                (vd-0) -- (gul)
                (vd-0) -- (gdr)
                (vd-1) -- (gur)
                (vd-1) -- (gdl);
                
                \draw[black, draw]
                (vl-0) -- (gur)
                (vl-0) -- (gdl)
                (vl-1) -- (gul)
                (vl-1) -- (gdl);
                \draw[draw= black, fill=white] (0,0) circle (0.25cm);
        \end{scope}
}

\newcommand{\drawGadgetBorder}[3]{
        \begin{scope}[name prefix=#1, shift={#2}, rotate=#3, scale=\gadgetScale]
                \drawEdgeVertexPair{vu}{rotate=90}
                \drawEdgeVertexPair{vr}{rotate=0}
                \drawEdgeVertexPair{vl}{rotate=-180}
                \drawGadgetVertices{}
                
                \draw[black, draw]
                (vu-0) -- (gul)
                (vu-0) -- (gdr)
                (vu-1) -- (gur)
                (vu-1) -- (gdl);
                
                \draw[black, draw]
                (vr-0) -- (gur)
                (vr-0) -- (gdl)
                (vr-1) -- (gul)
                (vr-1) -- (gdr);
                
                \draw[black, draw]
                (vl-0) -- (gur)
                (vl-0) -- (gdl)
                (vl-1) -- (gul)
                (vl-1) -- (gdl);
                \draw[draw= black, fill=white] (0,0) circle (0.25cm);
        \end{scope}
}
\newcommand{\drawGadgetCorner}[3]{
        \begin{scope}[name prefix=#1, shift={#2}, rotate=#3, scale=\gadgetScale]
                \drawEdgeVertexPair{vu}{rotate=90}
                \drawEdgeVertexPair{vr}{rotate=0}
                \drawGadgetVertices{}
                
                \draw[black, draw]
                (vu-0) -- (gul)
                (vu-0) -- (gdr)
                (vu-1) -- (gur)
                (vu-1) -- (gdl);
                
                \draw[black, draw]
                (vr-0) -- (gur)
                (vr-0) -- (gdl)
                (vr-1) -- (gul)
                (vr-1) -- (gdr);
                
                \draw[draw= black, fill=white] (0,0) circle (0.25cm);
        \end{scope}
}

\NewDocumentCommand{\drawConnection}{ O{black} O{red} m m m }{
        \ifthenelse{#5=0}{
                \draw[#1, draw, thick]
                (#3-0) -- (#4-1)
                (#3-1) -- (#4-0);
        }
        {
                \draw[#2, draw, thick]
                (#3-0) -- (#4-0)
                (#3-1) -- (#4-1);
        }
}

\newcommand{\drawBaseGraph}[4]{
        \def\length{7}
        \def\height{1}
        \begin{scope}[scale = 2, shift = {(-2,1)}, yscale=-1]
                \foreach \i in {0,...,\length} {
                        \foreach \j in {0,...,\height} {
                                \node[baseVertex] (u-\i-\j) at (\i,\j) {};
                        }
                }
                
                \pgfmathtruncatemacro{\prevlen}{\length-1}
                \pgfmathtruncatemacro{\prevhei}{\height-1}
                \foreach \i in {0,...,\prevlen} {
                        \pgfmathtruncatemacro{\nexti}{1+\i};
                        \foreach \j in {0,...,\height} {
                                \ifthenelse{\equal{\i}{#1} \AND \equal{\j}{#2} \AND \equal{\nexti}{#3} \AND \equal{\j}{#4}}
                                {
                                        \path[twistedBaseEdge] (u-\i-\j) to node{\includegraphics[width=6mm]{robber.png}} (u-\nexti-\j);
                                }
                                {
                                        \draw[baseEdge] (u-\i-\j) -- (u-\nexti-\j);
                                }
                        }
                }
                \foreach \i in {0,...,\length} {
                        \foreach \j in {0,...,\prevhei} {
                                \pgfmathtruncatemacro{\nextj}{1+\j};
                                \ifthenelse{\equal{\i}{#1} \AND \equal{\j}{#2} \AND \equal{\i}{#3} \AND \equal{\nextj}{#4}}
                                {\path[twistedBaseEdge] (u-\i-\j) to node{\includegraphics[width=6mm]{robber.png}} (u-\i-\nextj);}
                                {\draw[baseEdge] (u-\i-\j) -- (u-\i-\nextj);}
                        }
                }
                
        \end{scope}
}

\newcommand{\drawCFIgraph}[5]{
        \drawGadgetCorner{g-0-0}{(-4,2)}{-90}
        \drawGadgetBorder{b-0-1}{(-2,2)}{180}
        \drawGadgetBorder{b-0-2}{(-0,2)}{180}
        \drawGadgetBorder{b-0-3}{(2,2)}{180}
        \drawGadgetBorder{b-0-4}{(4,2)}{180}
        \drawGadgetCorner{g-0-5}{(6,2)}{180}
        
        \drawGadgetBorder{b-1-0}{(-4,0)}{-90}
        \drawGadgetCenter{c-1-1}{(-2,0)}
        \drawGadgetCenter{c-1-2}{(0,0)}
        \drawGadgetCenter{c-1-3}{(2,0)}
        \drawGadgetCenter{c-1-4}{(4,0)}
        \drawGadgetBorder{b-1-5}{(6,0)}{90}
        
        \drawGadgetCorner{g-2-0}{(-4,-2)}{0}
        \drawGadgetBorder{b-2-1}{(-2,-2)}{0}
        \drawGadgetBorder{b-2-2}{(-0,-2)}{0}
        \drawGadgetBorder{b-2-3}{(2,-2)}{0}
        \drawGadgetBorder{b-2-4}{(4,-2)}{0}
        \drawGadgetCorner{g-2-5}{(6,-2)}{90}
        
        \drawConnection{g-0-0vu}{b-0-1vr}{0}
        \drawConnection{b-0-1vl}{b-0-2vr}{0}
        \drawConnection[black][red]{b-0-2vl}{b-0-3vr}{#3}
        \drawConnection{b-0-3vl}{b-0-4vr}{0}
        \drawConnection{b-0-4vl}{g-0-5vr}{0}
        
        \drawConnection{g-0-0vr}{b-1-0vl}{0}
        \drawConnection{b-0-1vu}{c-1-1vu}{0}
        \drawConnection{b-0-2vu}{c-1-2vu}{0}
        \drawConnection[black][red]{b-0-3vu}{c-1-3vu}{#2}
        \drawConnection{b-0-4vu}{c-1-4vu}{0}
        \drawConnection{g-0-5vu}{b-1-5vr}{0}
        
        \drawConnection{b-1-0vu}{c-1-1vl}{0}
        \drawConnection{c-1-1vr}{c-1-2vl}{0}
        \drawConnection[black][red]{c-1-2vr}{c-1-3vl}{#1}
        \drawConnection{c-1-3vr}{c-1-4vl}{0}
        \drawConnection{c-1-4vr}{b-1-5vu}{0}
        
        \drawConnection{g-2-0vu}{b-1-0vr}{0}
        \drawConnection[black][red]{b-2-1vu}{c-1-1vd}{#5}
        \drawConnection{b-2-2vu}{c-1-2vd}{0}
        \drawConnection{b-2-3vu}{c-1-3vd}{0}
        \drawConnection{b-2-4vu}{c-1-4vd}{0}
        \drawConnection{g-2-5vr}{b-1-5vl}{0}
        
        \drawConnection[black][red]{g-2-0vr}{b-2-1vl}{#4}
        \drawConnection{b-2-1vr}{b-2-2vl}{0}
        \drawConnection{b-2-2vr}{b-2-3vl}{0}
        \drawConnection{b-2-3vr}{b-2-4vl}{0}
        \drawConnection{b-2-4vr}{g-2-5vu}{0}
}

\newcommand{\drawPebble}[2]{\draw[draw=#1!70!black, fill=#1] (#2) circle (1.5mm);}

\newcommand{\highlightSubgraph}[4]{
        \begin{scope}[on behind layer]
                \draw[line width=1mm, draw=magenta!50!white]
                (#1) -- (#2)
                (#1) -- (#3)
                (#1) -- (#4)
                (#2) -- (#3)
                (#2) -- (#4)
                (#3) -- (#4);
        \end{scope}
}
\tikzstyle{edgeVertex} = [circle, fill=black, inner sep=0, minimum size =1mm]
\tikzstyle{baseVertex} = [basevertex, gray]
\tikzstyle{baseEdge} = [draw, gray, ultra thick]
\tikzstyle{twistedBaseEdge} = [draw, red, ultra thick]
%


\newcommand{\id}[1]{#1}

%
%
%
%
\newcommand{\drawgadget}[5] {
        
        \draw[dashed, thick, gray] (#3,#4) circle (0.65);
        
        \ifnum #1=1
        \node[vertex,#5] (#2\id{0}) at (#3,#4){};
        \else \ifnum #1=2
        \node[vertex,#5] (#2\id{11}) at (#3+0.25,#4-0.25){};
        \node[vertex,#5] (#2\id{00}) at (#3-0.25,#4+0.25){};
        
        \else
        \foreach \x in {0,...,1}
        \foreach \y in {0,...,1}
        \foreach \z in {0,...,1}
        {
                \ifodd \numexpr \x + \y + \z \else
                \node[vertex, #5] (#2\x\y\z) at ({#3 + 0.5*( \z - 0.3*\x) -0.2}, {#4 + 0.5*(\x + 0.7*\z - 1.7*\y -\x*\z)+0.1}){};
                \fi
        }
        \fi
        \fi
        
}

\newcounter{memcount}
\newcommand{\countmem}[2]{%
        \setcounter{memcount}{0}%
        \foreach \i in #1{%
                \ifthenelse{\equal{\i}{#2}} {%
                        \stepcounter{memcount}%
                }{}%
        }%
        
}


\newcommand{\suffx}{x}
\newcommand{\suffy}{y}
\newcommand{\suffz}{z}
\newcommand{\suffsum}{sum}
\newcommand{\suffseq}{seq}
\newcommand{\suffrr}{rr}

\newcommand{\foreachdim}[4] {
        \ifnum #1=1
        \expandafter\foreach \csname#2\suffx\endcsname in {0,...,1} {
                \expandafter\def\csname#2\suffsum\endcsname{%
                        \csname#2\suffx\endcsname}
                \expandafter\def\csname#2\suffseq\endcsname{%
                        \csname#2\suffx\endcsname}
                \expandafter\def\csname#2\suffrr\endcsname{%
                        {\csname#2\suffx\endcsname}}
                #3
        }
        \else \ifnum #1=2
        \expandafter\foreach \csname#2\suffx\endcsname in {0,...,1} {
                \expandafter\foreach \csname#2\suffy\endcsname in {0,...,1} {
                        \expandafter\def\csname#2\suffsum\endcsname{%
                                \csname#2\suffx\endcsname + %
                                \csname#2\suffy\endcsname}
                        \expandafter\def\csname#2\suffseq\endcsname{%
                                \csname#2\suffx\endcsname\csname#2\suffy\endcsname}
                        \expandafter\def\csname#2\suffrr\endcsname{%
                                {\csname#2\suffx\endcsname,\csname#2\suffy\endcsname}}
                        #3
                }
        }
        \else
        \expandafter\foreach \csname#2\suffx\endcsname in {0,...,1} {
                \expandafter\foreach \csname#2\suffy\endcsname in {0,...,1} {
                        \expandafter\foreach \csname#2\suffz\endcsname in {0,...,1} {
                                \expandafter\def\csname#2\suffsum\endcsname{%
                                        \csname#2\suffx\endcsname + %
                                        \csname#2\suffy\endcsname +%
                                        \csname#2\suffz\endcsname}
                                \expandafter\def\csname#2\suffseq\endcsname{%
                                        \csname#2\suffx\endcsname\csname#2\suffy\endcsname\csname#2\suffz\endcsname}
                                \expandafter\def\csname#2\suffrr\endcsname{%
                                        {\csname#2\suffx\endcsname,\csname#2\suffy\endcsname,\csname#2\suffz\endcsname}}
                                #3
                        }
                }
        }
        \fi\fi
        
}

%
%
%
%
%
%
%
%
%
\newcommand{\drawconnectgadget}[9] {
        
        \foreachdim{#1}{a} {
                \foreachdim{#2}{b}{
                        \pgfmathsetmacro\as{\arr[#5]}
                        \pgfmathsetmacro\bs{\brr[#6]}
                        \ifodd \numexpr \asum \else
                        \ifodd \numexpr \bsum \else
                        \ifthenelse{
                                \(\(\NOT \equal{#9}{twist}\) \AND \equal{\as}{\bs}\) \OR %
                                \(\equal{#9}{twist} \AND \(\NOT \equal{\as}{\bs}\)\)
                        }{
                                \countmem{{#7}}{\aseq\id{E}\bseq}
                                \ifthenelse{\value{memcount} > 0}{
                                        \path (#3\aseq) edge[bend right=5] (#4\bseq);
                                }{
                                        \countmem{{#8}}{\aseq\id{E}\bseq}
                                        \ifthenelse{\value{memcount} > 0}{
                                                \path (#3\aseq) edge[bend left=5] (#4\bseq);
                                        }{
                                                \path (#3\aseq) edge (#4\bseq);
                                        }
                                }
                        }{}
                        \fi\fi
                }
        }
        
}

%
%
%
\newcommand{\gadgetlabel}[5] {
        \node[fill = none] () at ($(#1#2#3#4) + #5$) {\scriptsize$#1(#2,#3,#4)$};
}



\tikzstyle{basevertex} = [circle, fill=black, inner sep=0mm, minimum size = 6mm]

\colorlet{gcolA}{colH}
\colorlet{gcolB}{colB}
\colorlet{gcolC}{colC}
\colorlet{gcolD}{colD}
\colorlet{gcolE}{colF}
\colorlet{gcolF}{colE}
\colorlet{gcolG}{colG}
\colorlet{gcolH}{colA}

 %
%

\usepackage{fontspec,thmtools,thm-restate,ifthen}
\usepackage[margin=22mm]{geometry}
\newtheorem{theorem}{Theorem}
\newtheorem{lemma}[theorem]{Lemma}
\newtheorem{corollary}[theorem]{Corollary}
\newtheorem{proposition}[theorem]{Proposition}
\newtheorem{definition}[theorem]{Definition}
\newtheorem{remark}[theorem]{Remark}
\newtheorem{claim}[theorem]{Claim}
\newtheorem{example}[theorem]{Example}
\newcommand{\confORfull}[2]{#1}
\newcommand{\ifAppendix}[1]{#1}

\newenvironment{claimproof}{\begin{proof}[Proof of claim]}{\end{proof}}

\newcommand{\claimqedhere}{\qedhere}
